% sec-02: the ten recipients, their mechanisms, and their asks.  Figure 3 (d2).
\section{The Ten Applications}

\draftinstr{full-apply writes 5,000 to 6,000 characters from
\appfile{funding/pdac-funding-applications/applications/README.md} \S1, \S2 and
\S7, plus each of the ten \appfile{applications/app-*/README.md} files for the
per-recipient ask and anchor. Do not restate the applications; state what the
set as a whole demonstrates.}

\subsection{Selection Rule}

\draftinstr{One paragraph: a recipient qualifies only if the report names its
program or the mechanism it runs. Give the anchor chunk for each of the ten and
state that no recipient was included on plausibility alone.}

\figslot{d2-type / scored grid}{6.6}{Ten applications on four attributes. The
science column is constant by construction: the same trial is described ten
times, and only the mechanism, the term, and the ask change.}

\subsection{What Differs, and What Does Not}

\draftinstr{full-apply builds Table 2: a full-width table with columns Number,
Recipient, Mechanism family, Term, Direct ask, and Perspective. Ten rows. Draw
values from the application READMEs; do not recompute them.}

\draftinstr{Then two paragraphs on the deliberate constant: the trial, the
population, the endpoints and the governance are identical across all ten, so a
reviewer comparing two of them sees only a mechanism difference.}

\subsection{Redundancy}

\figslot{graphviz-type / dependency DAG}{6.2}{Ten awards against eight trial
activities. Six activities have two or more upstream awards and survive any
single refusal; the two that do not are drawn heavier, and both are regulatory.}

\draftinstr{One paragraph reading the figure: name the two single-source
activities, IND activation and IRB approval, and say why they are funded first
in every application that can fund them.}
